\subsection*{The Host Program}

The host side is plain C++ with two Croqtile API calls: one creates a shaped buffer from raw storage, the other hands a view of it to the kernel.

\begin{lstlisting}[language=croqtile,caption={Creating the shaped inputs}]
auto lhs = croq::make_spandata<croq::s32>(4, 8);
auto rhs = croq::make_spandata<croq::s32>(4, 8);
lhs.fill_random(-10, 10);
rhs.fill_random(-10, 10);
\end{lstlisting}

\texttt{croq:\allowbreak{}:\allowbreak{}make\_\allowbreak{}spandata\textless{}T\textgreater{}(dims.\allowbreak{}.\allowbreak{}.\allowbreak{})} creates an owning tensor buffer on the host. You pass the element type as a template parameter and the dimensions as arguments. \texttt{fill\_random} populates it with values in the given range.

\begin{lstlisting}[language=croqtile,caption={Calling the kernel}]
auto res = ele_add(lhs.view(), rhs.view());
\end{lstlisting}

Calling a \texttt{\_\_co\_\_} function from C++ looks like a normal function call. The \texttt{.view()} method produces a non-owning \texttt{spanned\_view} from an owning \texttt{spanned\_data} --- this is how you pass host tensors into a Croqtile function without transferring ownership. The return value \texttt{res} is a \texttt{croq:\allowbreak{}:\allowbreak{}spanned\_\allowbreak{}data} that owns its buffer.

The rest of \texttt{main} is ordinary verification: loop over every element and assert equality. The host program is deliberately boring --- the interesting work happens in the \texttt{\_\_co\_\_} function.
